% Pursuing Stacks (cotes 134-2 to 134-8): archive cote 134-4, site page 63 (PDF page 64), typescript page 151.
% typescript. Draft compared: ps3.tex:3201-3278.
% AI-transcribed by Claude Opus 5.5 (claude-opus-5-5, adaptive thinking, effort high), Message Batches API, batch ps2 (msgbatch_01NVwr1Yuo8qJf5rykeGBXPK), 2026-11-12; not yet checked by eye.
\scanpage{134-4/63}
\mspage{151}

The condition a) will be tacitly understood throughout, when taking
a $\underline{W}$ to replace usual weak equivalences. Conditions b) and c) \inserted{then} were seen to
be enough to imply that $h_{\underline{W}} = h_{(Cat)}$ . On the other hand, one sees at
once that for the proposition above for a \deleted{topos A} topos $A\hat{}$ \deleted{of the typ}\hand{wh}ich
is strictly totally aspheric, when we define now $\underline{W}_A \subset \mathrm{F}\over{l}{\ell}(A\hand{\hat{}}\ )$ as just
$i_A^{-1}(\underline{W})$ , in order to conclude $h_{(\underline{W}_A)} = h_{A\hat{}}$ , all we made use of was
(besides saturation of $\underline{W}$ of course i.e. a)) b) and d).

One may object that d) is'nt expressed in terms of the category
structure of (Cat) only, but we could express it in terms of this struc-%
ture, by the remark that C has a final object iff there exists a\xout{n}
$\hand{\Delta}^1$-homotopy of $id_C$ to a constant section of C (this ``section'' will indeed
be defined necessarily by a final object of C). As was to be expected,
in this formulation, as in c) too, the object $\hand{\Delta}_1$ of (Cat) is playing a
crucial role. But at this point i\hand{t} occurs to me that c) implies d),
by the homotopy lemma - thus for the time being all we needed was
a)b)c).

6\over{\unsure{0}}{1}. We now in the long last get back to the ``key result'' promised time
ago, and which we kept pushing off. To pay off the trouble of the long
digression in between, maybe it'll come out more smoothly. It shall be
concerned with a functor
\[ (1) \qquad i : A \longrightarrow (Cat) \quad , \]
where in the end A will be (or turn out to be) a strict test category,
and we want to give characterizations for i to be a \inserted{(weak)} test functor, namely
the corresponding functor
\[ (2) \qquad i^{*} : (Cat) \longrightarrow A\hat{} \]
to induce an equivalence between the localizations, with respect to
``weak equivalences''.\hand{*} We'll now be a little more demanding, and instead \margin{* Plus a little more, see below (8).}
of just assuming it is the usual notion of weak equivalence either in (Cat\hand{)}
or in $\over{\illegible{1}}{A}\hat{}$, I'll assume that the set $\underline{W}_A \subset \mathrm{F1}(A\hat{})$ is defined in terms of
a saturated \xout{sub}set $\underline{W} \subset \mathrm{F}\over{l}{\ell}((Cat))$ of arrows in \hand{(}Cat) , by taking the invers
% NOTE: the right edge of the scan cuts the end of this line ("invers")
image by the functor
\[ (3) \qquad i_A : A\hat{} \rightarrow (Cat) \]
as above, whereas in the target category (Cat) of (2), we'll work
with another saturated set $\underline{W}'$ - thus besides \over{\illegible{1}}{(}1), the data are more-%
over
\[ (4) \qquad \underline{W} \ , \ \underline{W}' \subset \mathrm{F}\over{l}{\ell}(Cat)) \qquad , \]

\margin{63}
